
\documentclass[runningheads]{llncs}

\usepackage[T1]{fontenc}
\usepackage{graphicx}
\usepackage{array}
\usepackage{ragged2e}
\usepackage{braket}
\usepackage{amsmath}


\begin{document}
%
\title{Auto Quantum Machine Learning for Multisource Classification}

\author{First Author\inst{1}\orcidID{0000-1111-2222-3333} \and
Second Author\inst{2,3}\orcidID{1111-2222-3333-4444} \and
Third Author\inst{3}\orcidID{2222--3333-4444-5555}}
%
\authorrunning{F. Author et al.}
% First names are abbreviated in the running head.
% If there are more than two authors, 'et al.' is used.
%
\institute{Princeton University, Princeton NJ 08544, USA \and
Springer Heidelberg, Tiergartenstr. 17, 69121 Heidelberg, Germany
\email{lncs@springer.com}\\
\url{http://www.springer.com/gp/computer-science/lncs} \and
ABC Institute, Rupert-Karls-University Heidelberg, Heidelberg, Germany\\
\email{\{abc,lncs\}@uni-heidelberg.de}}
%
\maketitle              % typeset the header of the contribution
%
\begin{abstract}
The abstract should briefly summarize the contents of the paper in
150--250 words.

\keywords{First keyword  \and Second keyword \and Another keyword.}
\end{abstract}
%
%
%%%%%%%%%%%%%%%%%%%%%%%%%%%%%%%%%%%%%%%%%%%%%%%%%%%%%%%%%%%%%%%%%%%%%%%%%%
\section{Introduction}



%%%%%%%%%%%%%%%%%%%%%%%%%%%%%%%%%%%%%%%%%%%%%%%%%%%%%%%%%%%%%%%%%%%%%%%%%%
\section{Approach}

\begin{equation}
E(\sigma_z) = \bra{\phi(x)}U^\dagger M U \ket{\phi(x)}
\end{equation}

\begin{equation}
M = \sigma_z \otimes \sigma_z \otimes \sigma_z \otimes I \otimes I \otimes ...
\end{equation}

\begin{equation}
\ket{\phi(x)} = (R_x(x_0) \otimes R_x(x_1) \otimes ... \otimes R_x(x_n)) \ket{0} 
\end{equation}

\begin{equation}
U(\boldsymbol{\theta}) = \prod_{\ell=1}^{L} U_{\mathrm{BEL}}(\boldsymbol{\theta}^{(\ell)})
\end{equation}

\begin{equation}
U_{\mathrm{BEL}}(\boldsymbol{\theta})
=
\mathrm{CNOT}_{n,1}
\left(
\prod_{i=1}^{n-1} \mathrm{CNOT}_{i,i+1}
\right)
\left(
\bigotimes_{i=1}^{n} R_Y(\theta_i)
\right)
\end{equation}

\begin{align*}
R_Y(\theta) &= e^{-i \theta Y / 2}\\
&= \begin{pmatrix}
cos(\frac{\theta}{2}) & -sin(\frac{\theta}{2})  \\
sin(\frac{\theta}{2})  & cos(\frac{\theta}{2}) 
\end{pmatrix}
\end{align*}

\begin{equation}
\mathrm{CNOT}_{i,i+1}
=
\ket{0}\!\bra{0}_i \otimes I_{i+1}
+
\ket{1}\!\bra{1}_i \otimes X_{i+1}
\end{equation}



%%%%%%%%%%%%%%%%%%%%%%%%%%%%%%%%%%%%%%%%%%%%%%%%%%%%%%%%%%%%%%%%%%%%%%%%%%
\section{Results}

% ===== Hyperparameters: QMLP\_initial\_run =====
\begin{table}[t]
\centering
\caption{Hyperparameters for experiment QMLP\_initial\_run}\label{tab1}
\small
\setlength{\tabcolsep}{4pt}
\renewcommand{\arraystretch}{1.1}
\begin{tabular}{>{\RaggedRight\arraybackslash}p{0.35\linewidth}
                >{\RaggedRight\arraybackslash}p{0.55\linewidth}}
\hline
Parameter & Value \\
\hline
accelerator & \texttt{auto} \\
batch\_size & \texttt{32} \\
callbacks & \texttt{[<src.logging\_utils.EpochMetricsTracker object at 0x7021fc10c080>]} \\
devices & \texttt{auto} \\
digits & \texttt{[5, 6, 7]} \\
enable\_checkpointing & \texttt{False} \\
enable\_progress\_bar & \texttt{True} \\
hidden\_dim & \texttt{[128]} \\
input\_dim & \texttt{14} \\
logger & \texttt{False} \\
loss\_fn & \texttt{CrossEntropyLoss()} \\
lr & \texttt{0.001} \\
max\_epochs & \texttt{100} \\
n\_folds & \texttt{5} \\
n\_layers & \texttt{3} \\
n\_qubits & \texttt{6} \\
num\_classes & \texttt{3} \\
num\_sanity\_val\_steps & \texttt{0} \\
num\_workers & \texttt{4} \\
parent\_run\_name & \texttt{QMLP\_initial\_run} \\
seed & \texttt{42} \\
\hline
\end{tabular}
\end{table}

% ===== Final Metrics: QMLP\_initial\_run =====
\begin{table}[t]
\centering
\caption{Final metrics for QMLP\_initial\_run}\label{tab2}
\begin{tabular}{|l|c|}
\hline
Metric & Value \\
\hline
accuracy\_avg\_mean & 0.870 \\
accuracy\_avg\_std & 0.031 \\
accuracy\_mean & 0.872 \\
accuracy\_std & 0.030 \\
f1\_macro\_mean & 0.870 \\
f1\_macro\_std & 0.031 \\
f1\_micro\_mean & 0.872 \\
f1\_micro\_std & 0.030 \\
f1\_weighted\_mean & 0.872 \\
f1\_weighted\_std & 0.030 \\
precision\_macro\_mean & 0.873 \\
precision\_macro\_std & 0.031 \\
precision\_micro\_mean & 0.872 \\
precision\_micro\_std & 0.030 \\
precision\_weighted\_mean & 0.874 \\
precision\_weighted\_std & 0.030 \\
recall\_macro\_mean & 0.870 \\
recall\_macro\_std & 0.031 \\
recall\_micro\_mean & 0.872 \\
recall\_micro\_std & 0.030 \\
recall\_weighted\_mean & 0.872 \\
recall\_weighted\_std & 0.030 \\
\hline
\end{tabular}
\end{table}

% ===== CV Metrics per Fold: QMLP\_initial\_run =====
\begin{table}[t]
\centering
\caption{Cross-validation metrics per fold for QMLP\_initial\_run}\label{tab3}
\begin{tabular}{|l|c|c|c|c|c|}
\hline
Metric & Fold 5 & Fold 4 & Fold 3 & Fold 2 & Fold 1 \\
\hline
accuracy & 0.861 & 0.890 & 0.865 & 0.912 & 0.833 \\
accuracy\_avg & 0.857 & 0.891 & 0.864 & 0.910 & 0.830 \\
f1\_macro & 0.858 & 0.890 & 0.864 & 0.911 & 0.829 \\
f1\_micro & 0.861 & 0.890 & 0.865 & 0.912 & 0.833 \\
f1\_weighted & 0.860 & 0.891 & 0.866 & 0.912 & 0.832 \\
precision\_macro & 0.864 & 0.890 & 0.864 & 0.912 & 0.831 \\
precision\_micro & 0.861 & 0.890 & 0.865 & 0.912 & 0.833 \\
precision\_weighted & 0.863 & 0.893 & 0.867 & 0.912 & 0.835 \\
recall\_macro & 0.857 & 0.891 & 0.864 & 0.910 & 0.830 \\
recall\_micro & 0.861 & 0.890 & 0.865 & 0.912 & 0.833 \\
recall\_weighted & 0.861 & 0.890 & 0.865 & 0.912 & 0.833 \\
train\_acc & 0.875 & 0.896 & 0.872 & 0.916 & 0.841 \\
train\_f1 & 0.877 & 0.895 & 0.871 & 0.917 & 0.839 \\
train\_loss & 0.630 & 0.476 & 0.637 & 0.577 & 0.661 \\
val\_acc & 0.858 & 0.889 & 0.864 & 0.911 & 0.829 \\
val\_f1 & 0.860 & 0.889 & 0.864 & 0.912 & 0.828 \\
val\_loss & 0.649 & 0.490 & 0.654 & 0.589 & 0.670 \\
\hline
\end{tabular}
\end{table}



%%%%%%%%%%%%%%%%%%%%%%%%%%%%%%%%%%%%%%%%%%%%%%%%%%%%%%%%%%%%%%%%%%%%%%%%%%
\section{Conclusion}







%%%%%%%%%%%%%%%%%%%%%%%%%%%%%%%%%%%%%%%%%%%%%%%%%%%%%%%%%%%%%%%%%%%%%%%%%%
\begin{credits}
\subsubsection{\ackname} A bold run-in heading in small font size at the end of the paper is
used for general acknowledgments, for example: This study was funded
by X (grant number Y).

\subsubsection{\discintname}
The authors have no competing interests to declare that are
relevant to the content of this article.
\end{credits}
%
% ---- Bibliography ----
%
% BibTeX users should specify bibliography style 'splncs04'.
% References will then be sorted and formatted in the correct style.
%
% \bibliographystyle{splncs04}
% \bibliography{mybibliography}
%
\begin{thebibliography}{8}
\bibitem{ref_article1}
Author, F.: Article title. Journal \textbf{2}(5), 99--110 (2016)

\bibitem{ref_lncs1}
Author, F., Author, S.: Title of a proceedings paper. In: Editor,
F., Editor, S. (eds.) CONFERENCE 2016, LNCS, vol. 9999, pp. 1--13.
Springer, Heidelberg (2016). \doi{10.10007/1234567890}


\bibitem{ref_url1}
LNCS Homepage, \url{http://www.springer.com/lncs}, last accessed 2023/10/25
\end{thebibliography}






\end{document}
